% §5 Experiments.

\section{Experiments}
\label{sec:exp}

We evaluate seven frontier models end-to-end on \bench{}, covering five proprietary models and two open-source models. This section describes the experimental setup (\S\ref{sec:exp-setup}) and presents the main leaderboard together with the per-scenario breakdown (\S\ref{sec:exp-main}). Turn-level trajectory analysis and failure taxonomy follow in \S\ref{sec:analysis}; case studies are deferred to Appendix~\ref{app:case-studies}.

\subsection{Experimental Setup}
\label{sec:exp-setup}
\label{sec:protocol}

\paragraph{Models.} We evaluate five proprietary models (\textbf{Claude Sonnet 4.6}, \textbf{Claude Opus 4.6}, \textbf{GPT-5.4 (high)}, \textbf{Gemini 3.1 Pro Preview}, and \textbf{Qwen 3.6 Plus}) and two open-source models (\textbf{Kimi K2.5} and \textbf{Kimi K2.6}). The Kimi K2.5 result reported here is from the public Infinigence-hosted endpoint; an internal instance-tuned Kimi K2.5 variant exists but is not reported in the main table to avoid mixing a public and a private baseline.

\paragraph{Framework and infrastructure.} Every model runs under a single agent framework, \textbf{OpenClaw}, with identical tool schemas across models; no per-model prompt engineering is performed. For the Kimi-series models, we apply the upstream fix for incorrect tool-call identifier sanitisation in OpenClaw\footnote{\url{https://github.com/openclaw/openclaw/issues/62319}}. Each task executes inside an isolated \texttt{docker-compose} group comprising the agent container, GreenMail for SMTP/IMAP, a Notion-compatible knowledge base, a Google-Sheets-compatible spreadsheet, and a Radicale CalDAV server. Containers are torn down between tasks, so runs do not share state.

\paragraph{Evaluation settings and metric.} All models use provider-default inference parameters, with \emph{extended thinking} enabled where supported (Claude, GPT-5.4, Gemini) and prompt caching enabled where providers offer it. Weighted score and Task Success follow \S\ref{sec:design-evaluation}; both are reported on a 0--100 scale throughout this section. We additionally report wall-clock time for a full benchmark sweep and total input/output tokens. \textbf{Main-table results are from a single full sweep per model.} We do not report run-to-run variance in this release; rankings among models that fall within a small weighted-score band should be read with that caveat in mind.

\paragraph{Cost normalisation.} For the efficiency view in \S\ref{sec:exp-main}, we treat total tool calls and total (input + output) tokens as compute-side proxies rather than as direct dollar-cost signals: API unit pricing differs by provider and changes month to month, so a frozen cost-per-token table would misrepresent at least one model by the time of camera-ready. Readers who need a monetary bound can recombine the per-model token counts in Table~\ref{tab:leaderboard} with the provider's published rate at read time.

\subsection{Main Results}
\label{sec:exp-main}

Table~\ref{tab:leaderboard} reports the overall leaderboard across seven models and 100 tasks. By weighted score, Claude Sonnet 4.6 (75.8), Claude Opus 4.6 (74.6), and GPT-5.4 (72.0) cluster within a 3.8\,pp band; because each model is evaluated with a single full sweep, rankings among models within this narrow band should be interpreted cautiously, and the larger gap to GPT-5.4 should likewise be read in the absence of repeated sweeps to quantify run-to-run variance. Under the stricter Task Success metric of Eq.~\eqref{eq:task-success}, the ordering changes: Claude Opus 4.6 leads at 20.0, followed by Sonnet 4.6 at 14.0 and GPT-5.4 at 9.0, and fully correct end-to-end completion is much rarer than partial progress.

% T4 --- Main leaderboard.
% Referenced from sections/05_experiments.tex.

\begin{table}[t]
  \caption{Main results on \bench{} (single-sweep). \emph{Score} is the mean of Eq.~\eqref{eq:score} across the 100 tasks, reported on a 0--100 scale. \emph{Task Success} follows Eq.~\eqref{eq:task-success}. \emph{Red-line fail} is the fraction of red-line checker evaluations the model fails, aggregated over the 55 red-line checkers (distributed across the 8 scenarios that carry any; see Table~\ref{tab:scenarios}). \emph{Wall time} is the total wall-clock for a full single-run sweep. \emph{Input} merges \texttt{input + cacheRead + cacheWrite} tokens so providers are comparable regardless of prompt-cache use. \emph{Tool calls} is the total across 100 tasks (mean 48--71 per task). Main-table results are from a single full sweep per model; run-to-run variance is not reported in this release, so ranking claims among models that fall within a small weighted-score band should be read as tentative.}
  \label{tab:leaderboard}
  \centering
  \scriptsize
  \setlength{\tabcolsep}{5pt}
  \begin{tabular}{lrrrrrrr}
    \toprule
    Model                         & Score & Task Success & Red-line fail & Wall time & Input tok.\ & Output tok.\ & Tool calls \\
    \midrule
    Claude Sonnet 4.6             & \textbf{75.8} & 14.0          &  3.6\% & 22.3\,h & 257.8\,M & 2.57\,M & 5{,}736 \\
    Claude Opus 4.6               & 74.6          & \textbf{20.0} &  5.5\% & 22.6\,h & 266.7\,M & 2.02\,M & 6{,}112 \\
    GPT-5.4 (high)                & 72.0          & 9.0           &  3.6\% & 26.1\,h & 231.5\,M & 2.93\,M & 7{,}052 \\
    Kimi K2.6                     & 68.4          & 7.0           &  7.3\% & 22.6\,h & 226.3\,M & 2.30\,M & 6{,}026 \\
    Gemini 3.1 Pro Preview        & 68.2          & 8.0           &  3.6\% & 18.9\,h & 338.8\,M & 1.77\,M & 5{,}877 \\
    Qwen 3.6 Plus                 & 57.2          & 5.0           & 14.5\% & 33.3\,h & 315.1\,M & 4.56\,M & 6{,}119 \\
    Kimi K2.5                     & 56.0          & 0.0           &  9.1\% & 22.8\,h & 214.0\,M & 1.47\,M & 4{,}776 \\
    \bottomrule
  \end{tabular}
\end{table}


\paragraph{Both metrics leave room for improvement.} No model exceeds 75.8 weighted score overall. Excluding the single-task EDA case, the highest per-scenario score is Claude Opus 4.6 at 92.6 on real estate, still leaving visible headroom. The stricter Task Success metric makes the gap clearer: even the strongest model fully solves only 20.0\% of tasks, Sonnet 4.6 fully solves 14.0\%, GPT-5.4 fully solves 9.0\%, and Kimi K2.5 fully solves 0.0\%. On the hardest scenario, project management, every model scores below 44.0, and the mean weighted score across seven models is about 35.1. A partially correct trajectory that misses a single silent mutation, an incomplete backend writeback, or a turn-specific rubric item still forfeits Task Success credit, which is consistent with \bench{}'s task design.

\paragraph{Per-scenario best is distributed across four models, not one.} Table~\ref{tab:scenario-heatmap} shows outright-best-by-scenario splits among four models: Claude Sonnet 4.6 on clinical assistant, e-commerce, HR, legal, and research assistant (five); Claude Opus 4.6 on content operation, insurance, journalist, project management, and real estate (five); GPT-5.4 on executive assistant; and Gemini 3.1 Pro Preview on investment analyst. The two Anthropic models tie on the single EDA task (100.0 each). Below third place the ranking reshuffles more sharply: Kimi K2.6 is competitive with Gemini on investment analyst (82.1 vs.\ 82.9), but trails it sharply on EDA (8.7 vs.\ 91.3) due to a single vision-dependent task it routed incorrectly. Coworker-agent evaluation therefore does not collapse to a single frontier-model ordering, and specialisation-driven scenarios are where models most separate: EDA separates Gemini from the Kimi family (with the caveat that EDA contains a single task and should be read as a case-level result rather than a stable scenario-level trend), project management is uniformly difficult for all models, and clinical / insurance / research-assistant are where red-line-heavy rubrics (Table~\ref{tab:scenarios}) reward compliance-aware trajectories.

\paragraph{No monotone relationship between score and tool or token consumption.} Three per-model point comparisons illustrate: +23\% tool calls at $-3.8$\,pp score (GPT-5.4 vs.\ Sonnet), +31\% input tokens at $-7.6$\,pp (Gemini vs.\ Sonnet), and $1.8\times$ output tokens at $-18.6$\,pp (Qwen vs.\ Sonnet). On \emph{score per thousand tool calls} (a compute-side proxy for action efficiency; see Cost normalisation in \S\ref{sec:exp-setup}), Sonnet 4.6 leads at 13.2, followed by Opus 4.6 (12.2), Kimi K2.5 (11.7), Gemini (11.6), Kimi K2.6 (11.3), GPT-5.4 (10.2), and Qwen (9.3). The two top-scoring models are also the two most tool-efficient, so under this proxy score and efficiency move together rather than trade off.

% T5 --- Per-scenario score heatmap (7 models x 13 scenarios).
% Referenced from sections/05_experiments.tex.
% Cell shading via \hm{N} macro (N = score * 100), defined in neurips_2026.tex.

\begin{table}[t]
  \caption{Per-scenario score on a 0--100 scale. Bold marks the per-scenario best. Cell shading is proportional to score (darker = higher). The per-scenario best is distributed across four models (Claude Sonnet 4.6, Claude Opus 4.6, GPT-5.4, and Gemini 3.1 Pro Preview); the two Anthropic models tie on EDA at 100.0. EDA contains a single task (Table~\ref{tab:scenarios}), so its row should be read as a case-level result rather than a stable scenario-level trend.}
  \label{tab:scenario-heatmap}
  \centering
  \footnotesize
  \setlength{\tabcolsep}{3.3pt}
  \begin{tabular}{lccccccc}
    \toprule
    Scenario                 & Sonnet 4.6 & Opus 4.6 & GPT-5.4 & Kimi K2.6 & Gemini 3.1 & Qwen 3.6 & Kimi K2.5 \\
    \midrule
    Clinical Assistant       & \hm{70}\textbf{92.3} & \hm{67}88.3 & \hm{66}86.9 & \hm{66}87.4 & \hm{55}73.2 & \hm{59}77.9 & \hm{53}71.0 \\
    Content Operation        & \hm{50}66.2 & \hm{58}\textbf{77.2} & \hm{54}71.2 & \hm{55}73.2 & \hm{55}72.1 & \hm{50}66.9 & \hm{41}54.9 \\
    E-commerce               & \hm{47}\textbf{61.8} & \hm{43}57.4 & \hm{38}50.3 & \hm{38}50.4 & \hm{42}55.5 & \hm{34}45.6 & \hm{35}46.3 \\
    EDA                      & \hm{80}\textbf{100.0} & \hm{80}\textbf{100.0} & \hm{72}95.7 & \hm{6}8.7 & \hm{68}91.3 & \hm{33}43.5 & \hm{33}43.5 \\
    Executive Assistant      & \hm{55}73.9 & \hm{57}75.6 & \hm{58}\textbf{76.7} & \hm{51}67.8 & \hm{37}49.2 & \hm{47}62.9 & \hm{42}56.0 \\
    HR                       & \hm{62}\textbf{82.5} & \hm{56}75.1 & \hm{61}81.0 & \hm{52}69.7 & \hm{58}77.5 & \hm{44}59.0 & \hm{50}66.1 \\
    Insurance                & \hm{62}81.9 & \hm{68}\textbf{90.3} & \hm{65}87.0 & \hm{57}75.9 & \hm{56}74.8 & \hm{47}62.5 & \hm{46}61.5 \\
    Investment Analyst       & \hm{61}80.8 & \hm{52}68.9 & \hm{62}82.0 & \hm{62}82.1 & \hm{62}\textbf{82.9} & \hm{23}30.9 & \hm{46}61.7 \\
    Journalist               & \hm{64}85.6 & \hm{65}\textbf{86.2} & \hm{60}79.8 & \hm{59}79.2 & \hm{56}74.7 & \hm{48}63.4 & \hm{43}57.1 \\
    Legal Assistant          & \hm{63}\textbf{84.1} & \hm{52}69.4 & \hm{45}59.5 & \hm{47}62.3 & \hm{54}72.2 & \hm{35}46.3 & \hm{28}37.0 \\
    Project Management       & \hm{31}41.7 & \hm{33}\textbf{43.6} & \hm{28}37.9 & \hm{20}26.9 & \hm{29}39.3 & \hm{17}22.6 & \hm{25}33.6 \\
    Real Estate              & \hm{67}89.1 & \hm{69}\textbf{92.6} & \hm{65}87.0 & \hm{63}84.5 & \hm{62}82.6 & \hm{66}87.9 & \hm{54}72.6 \\
    Research Assistant       & \hm{61}\textbf{81.4} & \hm{58}77.2 & \hm{55}73.1 & \hm{58}76.9 & \hm{51}67.9 & \hm{47}62.8 & \hm{44}59.2 \\
    \bottomrule
  \end{tabular}
\end{table}


\label{sec:exp-scenario}
